
\documentclass[11pt]{article}

\usepackage[margin=1in]{geometry}
\usepackage{amsmath,amssymb,amsthm}
\usepackage{listings}
\usepackage{xcolor}
\usepackage{graphicx}
\usepackage{enumitem}
\usepackage{hyperref}

\definecolor{codegreen}{rgb}{0,0.6,0}
\definecolor{codegray}{rgb}{0.5,0.5,0.5}
\definecolor{codepurple}{rgb}{0.58,0,0.82}
\definecolor{backcolour}{rgb}{0.95,0.95,0.92}

\lstdefinestyle{mystyle}{
    backgroundcolor=\color{backcolour},
    commentstyle=\color{codegreen},
    keywordstyle=\color{blue},
    numberstyle=\tiny\color{codegray},
    stringstyle=\color{codepurple},
    basicstyle=\ttfamily\small,
    breaklines=true,
    captionpos=b,
    keepspaces=true,
    numbers=left,
    numbersep=5pt,
    showspaces=false,
    showstringspaces=false,
    showtabs=false,
    tabsize=4
}

\lstset{style=mystyle}

\title{CPSC 250 \\ Lecture 1: Primitive Data Types and Memory}
\author{Edward Brash}
\date{}

\begin{document}

\maketitle

\section{Introduction}

A common misconception among students entering CPSC 250 is that they should already remember every detail from CPSC 150/150L perfectly. In reality, students enter with a wide range of preparation levels.

The goal of this course is not only to reinforce earlier material, but also to develop a much deeper understanding of how Python actually works internally.

In particular, we will focus heavily on:

\begin{itemize}
    \item Data structures
    \item Memory representation
    \item Algorithmic thinking
    \item Object-oriented programming
    \item Numerical data representation
\end{itemize}

\section{Why Data Structures Matter}

In many ways, data structures are the central idea in programming.

Every piece of data in a computer program must ultimately be stored somewhere in memory. Understanding:

\begin{itemize}
    \item how data is represented,
    \item how it is organized,
    \item and how efficiently it can be manipulated
\end{itemize}

is essential for designing good algorithms.

Python intentionally hides many low-level implementation details from programmers. This is convenient, but it can also obscure what is really happening underneath the hood.

In this course, we will repeatedly pull back that abstraction layer.

\section{Primitive Data Types in Python}

The four basic primitive-style data types we will review are:

\begin{enumerate}
    \item Integers
    \item Floating-point numbers
    \item Strings
    \item Booleans
\end{enumerate}

\subsection{Examples}

\begin{center}
\begin{tabular}{|c|c|}
\hline
Type & Examples \\
\hline
Integer & \texttt{1}, \texttt{0}, \texttt{42}, \texttt{-6} \\
Float & \texttt{1.2}, \texttt{-3.14159}, \texttt{3.28e-6} \\
String & \verb|"a"|, \verb|"Bob"|, \verb|\n| \\
Boolean & \texttt{True}, \texttt{False} \\
\hline
\end{tabular}
\end{center}

\subsection{Important Notes}

Different data types:

\begin{itemize}
    \item require different amounts of memory,
    \item use different storage mechanisms,
    \item and support different operations.
\end{itemize}

\section{Integers in Python}

One of Python's nicest features is that integers effectively behave as arbitrary-precision numbers.

Unlike C or C++, Python integers do not normally overflow when they exceed a fixed size.

\subsection{Overflow in C/C++}

In low-level languages like C/C++, integer variables are allocated a fixed number of bytes.

For example:

\begin{itemize}
    \item 4-byte integers can store only a limited range of values.
    \item Exceeding that range produces overflow.
\end{itemize}

Python hides this complexity from the programmer.

\section{Python Integers are Objects}

Although integers appear to be primitive values, Python integers are actually full Python objects implemented internally in C.

Conceptually:

\begin{itemize}
    \item a Python integer stores metadata,
    \item bookkeeping information,
    \item and dynamically allocated memory.
\end{itemize}

This flexibility comes at the cost of additional memory overhead.

\subsection{Factorial Example}

\begin{lstlisting}[language=Python]
def factorial(n):
    if n == 0 or n == 1:
        return 1

    return n * factorial(n - 1)

print(factorial(231))
\end{lstlisting}

Python can compute extremely large integers with hundreds of digits automatically.

Internally, Python allocates additional memory as necessary.

\subsection{Memory Cost}

Even very small integers require substantial overhead in Python.

A Python integer may require roughly 28 bytes or more of memory, depending on the system architecture.

This is because Python integers behave more like objects than primitive machine-level integers.

\section{Floating Point Numbers}

Most programming languages represent floating point numbers using binary fractions.

This introduces important limitations:

\begin{enumerate}[label=(\alph*)]
    \item Some decimal numbers cannot be represented exactly.
    \item Floating-point arithmetic introduces rounding error.
\end{enumerate}

\section{Scientific Notation Review}

Recall ordinary scientific notation:

\[
3352.28 = 3.35228 \times 10^3
\]

This representation separates a number into:

\begin{itemize}
    \item mantissa/significand
    \item base
    \item exponent
\end{itemize}

\section{Binary Fractions}

Decimal fractions use powers of 10:

\[
32.56 = 3(10^1) + 2(10^0) + 5(10^{-1}) + 6(10^{-2})
\]

Binary fractions instead use powers of 2.

For example:

\[
1011.011_2
\]

corresponds to:

\[
1(2^3) + 0(2^2) + 1(2^1) + 1(2^0)
+ 0(2^{-1}) + 1(2^{-2}) + 1(2^{-3})
\]

which equals:

\[
8 + 0 + 2 + 1 + 0 + 0.25 + 0.125 = 11.375
\]

\section{Binary Scientific Notation}

Floating point numbers are stored in normalized binary scientific notation.

Example:

\[
11.01_2 = 1.101_2 \times 2^1
\]

The normalized form always begins with:

\[
1.xxxxx
\]

This allows one bit to be implicitly assumed rather than explicitly stored.

\section{32-Bit Floating Point Representation}

IEEE-754 single precision floating point numbers are stored using:

\begin{itemize}
    \item 1 sign bit
    \item 8 exponent bits
    \item 23 mantissa bits
\end{itemize}

\subsection{Exponent Bias}

The exponent field uses a bias of 127.

Thus:

\[
\text{stored exponent} = E + 127
\]

Examples:

\begin{align*}
00000001_2 &\rightarrow 2^{-126} \\
11111110_2 &\rightarrow 2^{127}
\end{align*}

Special exponent values are reserved for:

\begin{itemize}
    \item infinity
    \item NaN (Not a Number)
    \item denormalized numbers
\end{itemize}

\section{Special Exponent Values}

\begin{itemize}
    \item exponent bits are all 1,
    \item and the mantissa is all 0,
\end{itemize}

then the number represents:

\[
+\infty
\]

or:

\[
-\infty
\]

depending on the sign bit.

Examples that may generate infinity include:

\begin{lstlisting}[language=Python]
1.0 / 0.0
\end{lstlisting}

or arithmetic overflow.

\subsubsection{NaN (Not a Number)}

If:

\begin{itemize}
    \item exponent bits are all 1,
    \item but the mantissa is not all 0,
\end{itemize}

then the result is NaN.

NaN represents undefined or invalid numerical operations.

Examples include:

\begin{lstlisting}[language=Python]
0.0 / 0.0
math.sqrt(-1)
\end{lstlisting}

NaN has unusual comparison behavior:

\begin{lstlisting}[language=Python]
float('nan') == float('nan')
\end{lstlisting}

returns:

\begin{lstlisting}
False
\end{lstlisting}

This behavior is intentional and follows the IEEE-754 standard.

\subsection{Summary}

\begin{center}
\begin{tabular}{|c|c|}
\hline
Exponent Field & Interpretation \\
\hline
00000000 & Denormalized Numbers / Zero \\
00000001 to 11111110 & Normal Floating Point Numbers \\
11111111 & Infinity or NaN \\
\hline
\end{tabular}
\end{center}

\section{Strings in Python}

Strings are significantly more complicated internally than they first appear.

Consider:

\begin{lstlisting}[language=Python]
s1 = "abcdef"
\end{lstlisting}

Conceptually, the characters are stored sequentially in memory:

\begin{center}
\begin{tabular}{|c|}
\hline
a \\
\hline
b \\
\hline
c \\
\hline
d \\
\hline
e \\
\hline
f \\
\hline
\end{tabular}
\end{center}

Each character is represented internally using character encoding schemes such as ASCII or Unicode.

\subsection{Memory Addresses}

\begin{lstlisting}[language=Python]
print(s1)
print(id(s1))
\end{lstlisting}

The function \texttt{id()} returns the memory identity of the object.

\section{Booleans}

In Python, booleans are actually a subclass of integers.

Specifically:

\[
\texttt{True} = 1
\]

and

\[
\texttt{False} = 0
\]

Internally, booleans therefore behave much like integers.

\section{Key Takeaways}

\begin{itemize}
    \item Python hides many implementation details.
    \item Data representation strongly affects algorithm behavior.
    \item Integers, floats, strings, and booleans are all implemented as objects.
    \item Floating point numbers are approximations.
    \item Understanding memory representation is essential for serious programming.
\end{itemize}

\end{document}

